\documentclass[11pt]{article}
\usepackage[a4paper,margin=25mm]{geometry}
\usepackage{amsmath,amssymb,amsthm,mathtools}
\usepackage[hidelinks]{hyperref}
\allowdisplaybreaks[2]
\setlength{\emergencystretch}{2em}
\newtheorem{theorem}{Theorem}
\newtheorem{lemma}[theorem]{Lemma}
\newtheorem{proposition}[theorem]{Proposition}
\title{An inverse-linear bounded-range bound for the $1/3$--$2/3$ conjecture}
\author{Leo Gao\\[-1pt]{\small AI-assisted research}}
\date{November 26, 2026 \quad$\cdot$\quad Research draft}
\begin{document}
\maketitle
\begin{abstract}
We give a sharp gap-insertion bound for a comparable pair in a uniformly random linear extension. Applied to a displacement-maximizing triple, it yields an additive improvement of order $1/D$ over the Brightwell--Felsner--Trotter constant for finite nonchains of incomparability range at most $D$. We also describe an exact rational prefix certificate for ruling out candidate counterexamples.
\end{abstract}

\section{Statement}

For a finite partially ordered set $P$, let $f$ be a uniformly random linear extension, interpreted as a rank map. Write $x\parallel y$ for incomparability and let
\[
 \delta(P)=\max_{x\ne y}\min\{\Pr[f(x)<f(y)],\Pr[f(y)<f(x)]\}.
\]
The range of $P$ is $\max_{x\in P}|\{y\in P:x\parallel y\}|$. The $1/3$--$2/3$ conjecture asks whether $\delta(P)\ge 1/3$ for every finite nonchain. Put $C=(5-\sqrt5)/10$ and $K=5+3\sqrt5$. Theorem~1.3 of Aires, Chan, Pak and Panova~\cite{ACPP} gives an additive improvement over $C$ of $(D+1)^{-3(D+1)}/12$ when the range is at most $D$.

\begin{theorem}
For every finite nonchain $P$ of range at most $D$, where $D\ge2$,
\[
\delta(P)\ge C+\min\left\{\frac{67}{242}-C,\frac{C}{K(D+1)+1}\right\}.
\]
\end{theorem}

The additive improvement in this estimate is of order $1/D$ as $D$ grows. The unrestricted conjecture remains open.

\section{Gap insertion}

\subsection{Lemma 1 --- Sharp gap-insertion inequality}

Let \(Q\) be a finite poset of range at most \(D\), and let \(a\prec_Q b\).

Suppose there exists at least one linear extension in which \(a\) and \(b\) are not consecutive.

Then

\[
\boxed{
\Pr_Q[f(b)-f(a)\geq2]
\geq\frac1{D+1}.
}
\]

Moreover, this dependence on \(D\) is optimal.

\paragraph{Proof.}

Partition the linear extensions of \(Q\) into two classes:

\begin{itemize}
\item \(\mathcal A\): extensions in which \(a,b\) are consecutive.
\item \(\mathcal B\): extensions in which at least one element separates them.

\end{itemize}
By assumption, \(\mathcal B\ne\varnothing\).

Define the set of eligible elements

\[
T=Q\setminus
\bigl(\downarrow a\;\cup\;\uparrow b\bigr),
\]

where the principal downset and upset include \(a\) and \(b\), respectively.

Every element occurring between \(a,b\) in any linear extension must belong to \(T\). Consequently, \(T\ne\varnothing\).

Take an arbitrary extension in \(\mathcal A\).

If an element of \(T\) occurs before \(a\), choose the last such element \(w\), remove it and insert it between \(a,b\).

Otherwise, choose the first element of \(T\) occurring after \(b\), remove it and insert it between \(a,b\).



The selected element crosses only elements incomparable with it, preserving the validity of the extension.



To see this, consider the first case. Every element crossed by \(w\) belongs to \(\downarrow a\). If \(w\) were required to precede a crossed element \(u\), transitivity would imply

\[
w\prec_Q u\preceq_Q a,
\]

contradicting \(w\in T\).

Therefore moving \(w\) to the right violates no order relation. The second case is the order-dual argument.

This defines a map

\[
\Phi:\mathcal A\longrightarrow\mathcal B.
\]

Every image of this map contains precisely one element between \(a,b\), so the identity of the moved element \(w\) can be recovered uniquely.

A possible preimage can be recovered by moving \(w\) back to its original position.

An element incomparable with at most \(D\) others has at most \(D+1\) possible ranks in any linear extension. One rank is its current position, leaving at most \(D\) possible preimages.

Hence

\[
|\mathcal A|\leq D|\mathcal B|.
\]

It follows that

\[
\Pr_Q[f(b)-f(a)\geq2]
=
\frac{|\mathcal B|}
{|\mathcal A|+|\mathcal B|}
\geq\frac1{D+1}.
\]

\[
\qquad\square
\]

\paragraph{Sharpness.}

Take a chain of \(D\) elements and one isolated element \(w\).

There are exactly \(D+1\) linear extensions, corresponding to the possible insertion positions of \(w\).

For any consecutive pair \(a,b\) in the chain, exactly one extension places \(w\) between them.

Thus

\[
\Pr[f(b)-f(a)\geq2]=\frac1{D+1}.
\]

The lemma is therefore sharp for every \(D\geq2\).

\subsection{Lemma 2 --- Conditional gap inequality}

Let \(a,b\) be incomparable elements of a finite poset \(P\) of range at most \(D\).

If some linear extension places \(a\) before \(b\), with an intervening element, then

\[
\boxed{
\Pr[f(b)-f(a)\geq2]
\geq
\frac{\Pr[f(a)<f(b)]}{D+1}.
}
\]

\textbf{Proof.} Add the relation \(a\prec b\) to \(P\) and take its transitive closure, obtaining a new poset \(Q\).

Adding relations cannot increase the range. The uniform distribution on the linear extensions of \(Q\) is precisely the conditional distribution of the linear extensions of \(P\), given \(f(a)<f(b)\).

Applying Lemma 1 gives the result. \(\square\)

This conditional formulation is the crucial improvement over the general-purpose rare-event estimate employed in the October 2026 paper.

\section{Proof of the bound}

The proof uses the triple analysis of Brightwell, Felsner and Trotter~\cite{BFT} and the refinements of Aires, Chan, Pak and Panova~\cite{ACPP}.



\subsection{Step 1 --- Boundary displacement}

Let

\[
h(v)=\mathbb E[f(v)]
\]

be the expected rank of an element.

List the elements in nondecreasing order of expected rank:

\[
h(v_1)\leq h(v_2)\leq\cdots\leq h(v_n).
\]

Define

\[
d_i=h(v_i)-i,
\qquad
M=\max_i|d_i|.
\]

We may assume the incomparability graph is connected. Otherwise, its connected components form an ordinal sum, and the comparison probabilities inside each component are preserved.

The two-element case is immediate.

Since the incomparability graph is connected and nontrivial, the poset has at least two minimal elements. Choose a minimal element \(u\ne v_1\).

Moving \(u\) to the beginning of any linear extension produces another valid linear extension. Each resulting extension has at most \(D+1\) preimages.

Consequently,

\[
\Pr[f(u)=1]\geq\frac1{D+1}.
\]

Whenever \(f(u)=1\), the rank of \(v_1\) is at least two. Therefore

\[
\boxed{
M\geq h(v_1)-1\geq\frac1{D+1}.
}
\]

This independently replaces the extremely small boundary-displacement estimate used in the existing proof.

\subsection{Step 2 --- Select a critical triple}

Lemma~4.1 of~\cite{ACPP} guarantees a triple of consecutive elements in expected-rank order,

\[
(x,y,z)=(v_m,v_{m+1},v_{m+2}),
\]

that does not form a chain and satisfies

\[
h(z)-h(x)\leq2,
\]

with

\[
\max\{|d_m|,|d_{m+2}|\}=M.
\]

Such a triple is called a BFT triple.

Set

\[
t=\delta(P)-C.
\]

The desired conclusion already holds if

\[
t\geq\frac{67}{242}-C.
\]

We therefore assume

\[
\delta(P)<\frac{67}{242}.
\]

The established BFT triple classification, including the Case C estimate in Section~2.1 of~\cite{ACPP}, then forces

\[
x\prec z,\qquad x\parallel y,\qquad y\parallel z.
\]

\subsection{Step 3 --- Control the probability of separated inversions}

Define

\[
p=\Pr[f(y)<f(x)],
\]

\[
p'=\Pr[f(z)<f(y)].
\]

The expected-rank ordering and the established height-probability inequality imply that, under our assumption,

\[
p,p'\leq\delta(P)=C+t.
\]

The proof of Lemma~2.4 of~\cite{ACPP} supplies a nonnegative quantity \(R\) satisfying

\[
\frac{p+p'}2\geq C+\frac R K
\]

and

\[
R\geq
\Pr[f(x)-f(y)\geq2]
+
\Pr[f(y)-f(z)\geq2].
\]

In particular,

\[
R\leq Kt,
\qquad
p,p'\geq C-t.
\]

Suppose that some extension realizes a separated inversion of \(x,y\).

By the new conditional gap inequality,

\[
R\geq
\Pr[f(x)-f(y)\geq2]
\geq\frac{p}{D+1}
\geq\frac{C-t}{D+1}.
\]

Combining the inequalities gives

\[
Kt\geq\frac{C-t}{D+1}.
\]

Rearranging,

\[
\boxed{
t\geq\frac{C}{K(D+1)+1}.
}
\]

The argument for a separated inversion of \(y,z\) is identical.

Thus the desired theorem holds whenever either kind of separated inversion is possible.

\subsection{Step 4 --- Resolve the remaining structural case}

The only remaining possibility is that every reversed occurrence of either pair is consecutive:

\[
f(y)<f(x)\Longrightarrow f(x)=f(y)+1,
\]

\[
f(z)<f(y)\Longrightarrow f(y)=f(z)+1.
\]

First, \(z\) covers \(x\). Otherwise, an element strictly between \(x,z\) in the partial order, together with the positive-probability event \(f(x)<f(y)<f(z)\), would force

\[
h(z)-h(x)>2,
\]

a contradiction.

The consecutive-inversion property also implies that for every other element \(w\),

\[
w\prec y\quad\text{or}\quad x\prec w,
\]

and

\[
w\prec z\quad\text{or}\quad y\prec w.
\]

Indeed, if neither alternative in the first assertion held, a linear extension containing \(y,w,x\) in that relative order would exist, contradicting the consecutive-inversion property. The second assertion follows similarly.

If \(w\) were incomparable with \(y\), these implications would yield

\[
x\prec w\prec z,
\]

contradicting the cover property.

Consequently, \(y\) is incomparable only with \(x,z\).

Write

\[
L=\{w:w\prec y\},
\qquad
U=\{w:y\prec w\}.
\]

The ground set decomposes as

\[
P=L\sqcup\{x,y,z\}\sqcup U,
\]

and

\[
L\cup\{x\}\prec U\cup\{z\}.
\]

We give the displacement calculation from Claims~4.3--4.4 of~\cite{ACPP} under the structure just derived. All elements of $L$ come before $x$ in the expected-rank order, and all elements of $U$ follow $z$, because $L\prec y\prec U$ and $x,y,z$ are consecutive in that order. Thus $m=|L|+1$.

Set

\[
L'=L\cup\{x\},\qquad U'=U\cup\{z\}.
\]

Let \(r\) be the probability that \(x\) is last in a uniform extension of \(L'\), and let \(s\) be the probability that \(z\) is first in a uniform extension of \(U'\).

Counting the possible insertion positions of \(y\) gives

\[
p=\frac r{1+r+s},
\qquad
p'=\frac s{1+r+s}.
\]

Define

\[
a=\max(p,p').
\]

Then

\[
r,s\leq\frac a{1-2a}.
\]

Let \(\alpha\) denote the expected number of elements of \(L\) following \(x\), and \(\beta\) the expected number of elements of \(U\) preceding \(z\).

The same insertion count gives

\[
\alpha\geq(1-p)(1-r),
\]

\[
\beta\geq(1-p')(1-s).
\]

Hence, defining

\[
V(a)=\frac{(1-a)(1-3a)}{1-2a},
\]

we obtain

\[
\alpha,\beta\geq V(a).
\]

The endpoint displacements satisfy

\[
d_m=p-\alpha,
\qquad
d_{m+2}=\beta-p'.
\]

Since the triple has expected span at most two,

\[
d_{m+2}\leq d_m.
\]

Therefore

\[
V(a)-a\leq d_{m+2}
\leq d_m\leq a-V(a).
\]

By our selection of the triple,

\[
M\leq a-V(a).
\]

The algebraic estimate

\[
a-V(a)
=
\frac{5a-5a^2-1}{1-2a}
\leq6(a-C)
\]

holds for \(C\leq a<1/3\), as established in Claim~4.4 of~\cite{ACPP}.

Thus

\[
\delta(P)-C\geq\frac M6
\geq\frac1{6(D+1)}.
\]

This is stronger than the required inverse-linear bound.

Combining all cases proves

\[
\boxed{
\delta(P)\geq
C+\min\left\{
\frac{67}{242}-C,\,
\frac{C}{K(D+1)+1}
\right\}.
}
\]

\[
\qquad\square
\]

\section{Exact prefix certificates}

For a reference linear extension $v_1,\ldots,v_n$ of a range-$D$ poset, every linear extension satisfies $|f(v_i)-i|\le D$. Indeed, the rank change is the difference between the counts of elements crossing $v_i$ in opposite directions, and only incomparable elements can cross it. At cut $k$, every possible prefix ideal $I$ therefore contains the first $k-D$ reference elements and is contained within the first $k+D$, with indices clipped to $[1,n]$.

Suppose $P$ has no balanced pair. Orient each pair in its strict majority direction, whose probability exceeds $2/3$. Three directed majority comparisons cannot form a cycle. In any total order at most two comparisons in a directed cycle on three elements hold, whereas the sum of their probabilities would exceed two. A tournament without a directed triangle is transitive. Its linear order extends $P$ and has reversal probability below $1/3$ for every pair. Use this as the reference order.

For each candidate ideal $I$ of size $k$ and incomparable reference pair $v_a,v_b$ with $a<b$, let $q_I(a,b)$ be the fraction of extensions of $I$ that reverse the pair when both elements lie in $I$. Take $q_I(a,b)=1$ when only $v_b$ is in $I$, and $q_I(a,b)=0$ otherwise. This lower-bounds the conditional reversal probability given $I$. The prefix under that conditioning is uniform because each of its linear extensions has exactly the same number of completions.

\begin{proposition}
Suppose nonnegative rational weights $w_{ab}$, with at least one positive weight, satisfy
\[
 3\sum_{a<b}w_{ab}q_I(a,b)\ \,\ge\ \sum_{a<b}w_{ab}
\]
for every candidate prefix ideal $I$. Then the reference order cannot be the majority order of a counterexample.
\end{proposition}
\begin{proof}
Average the inequality over the distribution of the prefix ideal in a uniform extension. Every unconditional reversal probability is strictly below $1/3$ in the majority order of a counterexample, a contradiction.
\end{proof}

For a prescribed finite prefix, a certificate involving pairs within that prefix applies uniformly to its range-$D$ continuations when all candidate ideals at the cut lie within the prefix. The fraction $q_I$ can be computed by exact extension counts. The criterion leaves open whether such certificates suffice across all bounded-range instances.

\section{Exact checks}
The accompanying \texttt{verify.py} enumerates naturally labelled posets and checks the gap inequality for comparable pairs in which separation occurs. The computation supports the insertion lemma. The general bound follows from the proof above.

\section*{Disclosure}
This manuscript was prepared with AI assistance. Independent proof review and priority assessment are pending.

\begin{thebibliography}{9}
\bibitem{ACPP} M. Aires, S. H. Chan, I. Pak and G. Panova, \emph{Breaking the infinite barrier in the $1/3$--$2/3$ conjecture}, arXiv:2609.30888v2 (2026). \href{https://arxiv.org/html/2609.30888v2}{Preprint.}
\bibitem{BFT} G. Brightwell, S. Felsner and W. Trotter, \emph{Balancing pairs and the cross product conjecture}, Order 12 (1995), 327--349. \href{https://trotter.math.gatech.edu/papers/97.pdf}{Article.}
\end{thebibliography}
\end{document}
